% =====================================================================
%  1bat-ud3-8.tex  —  HORA 8: El punt òptim del casal: vèrtex i extrems
%  (sense preàmbul: s'inclou des de main.tex)
% =====================================================================
\horatitol{Sessió 8 — El punt òptim del casal: vèrtex i extrems}%
{Estudi de cas (Activitat 3): finançar les colònies venent samarretes; preus d'equilibri i preu de benefici màxim.}

\subsection{Idea clau}

A la sessió anterior vam aprendre a trobar el vèrtex i els talls d'una paràbola.
Avui ho apliquem a un cas real (Activitat 3): un grup de joves vol \textbf{finançar
les colònies} venent samarretes. Un estudi de mercat indica que el benefici depèn
del \textbf{preu de venda} $x$ segons
\[
B(x)=-5x^2+150x-400 \quad(\text{en euros}).
\]
La paràbola va cap avall ($a=-5<0$), o sigui que té un \textbf{màxim}: existeix un
preu que dóna el millor benefici possible. La nostra feina és trobar-lo i \emph{decidir}.

\begin{idea}
\textbf{Dues preguntes clau d'un problema d'optimització:}
\begin{itemize}[leftmargin=1.4em, itemsep=2pt]
  \item \textbf{Punts de tall amb l'eix $X$} ($B(x)=0$): els preus on \emph{no es
        guanya ni es perd} (benefici zero, el negoci queda \textbf{equilibrat}).
        Entre aquests dos preus, el benefici és positiu.
  \item \textbf{Vèrtex} ($x_V=\dfrac{-b}{2a}$): el preu que dóna el \textbf{benefici
        màxim}. És el \emph{punt òptim}, la resposta a «a quin preu hem de vendre?».
  \end{itemize}
\textbf{Decidir amb dades:} un cop trobat el punt òptim, cal \emph{justificar} la
decisió: per què aquest preu i no un altre.
\end{idea}

\subsection{Exemples}

\begin{exemple}[la paràbola del cas, llegida sencera]
Per a $B(x)=-5x^2+150x-400$, aquests són els punts que ens interessen:

\begin{center}
\begin{tikzpicture}
\begin{axis}[udaxis, width=10cm, height=6cm,
    xlabel={\small preu samarreta ($x$, \euro)}, ylabel={\small benefici (\euro)},
    xmin=0, xmax=30, ymin=-450, ymax=850,
    xtick={0,5,10,15,20,25,30}, ytick={-400,0,400,725}]
  \addplot[udblue, domain=2:28, samples=80]{-5*x^2+150*x-400};
  \addplot[udgrey, dashed, domain=0:30, samples=2]{0};
  \addplot[udnavy, only marks, mark=*, mark size=1.8pt] coordinates {(15,725) (2.96,0) (27.04,0)};
  \node[udnavy, font=\scriptsize, anchor=south] at (axis cs:15,725) {vèrtex $(15,\,725)$};
  \node[udnavy, font=\scriptsize, anchor=north] at (axis cs:3,0) {$\approx 3$};
  \node[udnavy, font=\scriptsize, anchor=north] at (axis cs:27,0) {$\approx 27$};
\end{axis}
\end{tikzpicture}

\figcap{$B(x)=-5x^2+150x-400$: entre $\approx 3\eur$ i $\approx 27\eur$ hi ha benefici; el màxim és a $x=15\eur$.}
\end{center}

\noindent La corba talla l'eix horitzontal a prop de $3\eur$ i $27\eur$ (preus
d'equilibri) i fa cim a $x=15\eur$, amb un benefici de $725\eur$.
\end{exemple}

\begin{exemple}[comprovar un valor per donar-li sentit]
Comprovem que a $x=15$ el benefici és realment $725\eur$:
\[
B(15)=-5\cdot 15^2+150\cdot 15-400=-5\cdot 225+2250-400=-1125+2250-400=725.
\]
I a un preu petit, com $x=5$: $B(5)=-125+750-400=225\eur$ (positiu, però lluny del
màxim). A $x=2$: $B(2)=-20+300-400=-120\eur$ (pèrdua: massa barat).
\end{exemple}

\subsection{Exercicis resolts}

\begin{exresolt}[preus d'equilibri: punts de tall]
Troba els preus a què el benefici $B(x)=-5x^2+150x-400$ és zero (ni guany ni
pèrdua).
\solucio
Resolem $-5x^2+150x-400=0$. Dividim tota l'equació entre $-5$ per simplificar:
\[
x^2-30x+80=0.
\]
Apliquem la fórmula de segon grau ($a=1$, $b=-30$, $c=80$):
\[
x=\frac{30\pm\sqrt{(-30)^2-4\cdot 1\cdot 80}}{2}
 =\frac{30\pm\sqrt{900-320}}{2}
 =\frac{30\pm\sqrt{580}}{2}.
\]
Com que $\sqrt{580}\approx 24,08$:
\[
x\approx\frac{30+24,08}{2}\approx 27,0
\qquad\text{i}\qquad
x\approx\frac{30-24,08}{2}\approx 3,0.
\]
\resposta{El benefici és zero a $\approx 3\eur$ i $\approx 27\eur$; entre tots dos preus el benefici és positiu.}
\end{exresolt}

\begin{exresolt}[el preu de benefici màxim: vèrtex]
Troba a quin preu cal vendre les samarretes per obtenir el benefici màxim, i quin és
aquest benefici.
\solucio
El preu òptim és l'abscissa del vèrtex, amb $a=-5$ i $b=150$:
\[
x_V=\frac{-b}{2a}=\frac{-150}{2\cdot(-5)}=\frac{-150}{-10}=15.
\]
El benefici màxim és $B(15)$:
\[
B(15)=-5\cdot 225+150\cdot 15-400=-1125+2250-400=725.
\]
\resposta{Cal vendre a $15\eur$ la samarreta; el benefici màxim és de $725\eur$.}
\end{exresolt}

\subsection{Exercicis per fer}

\begin{exercicis}
\begin{exlist}
  \item En el cas de les samarretes, comprova calculant que el benefici a $x=10\eur$
        i a $x=20\eur$ és el mateix. Per què té sentit que dos preus diferents donin el
        mateix benefici? (Pista: pensa en la simetria de la paràbola respecte al vèrtex.)
  \item Una altra parella modelitza la venda de polseres amb $B(x)=-2x^2+20x-32$.
        \begin{enumerate}[label=\alph*), leftmargin=1.6em, topsep=2pt]
          \item Troba els preus d'equilibri (resol $B(x)=0$).
          \item Troba el preu de benefici màxim (vèrtex) i quin és aquest benefici.
        \end{enumerate}
  \item Per a $B(x)=-5x^2+150x-400$, calcula el benefici a $x=8\eur$ i a $x=22\eur$ i
        digues, en cada cas, si convé apujar o abaixar el preu per acostar-se al màxim.
  \item \textbf{(Decisió justificada)} Amb les dades del cas de les samarretes, una
        parella proposa vendre a $25\eur$ «perquè com més car, més es guanya». Calcula
        $B(25)$ i explica, amb el resultat, per què aquesta idea és equivocada.
  \item \textbf{(Interpretació del signe)} Per què la paràbola del benefici va cap
        avall i no cap amunt? Què passaria, en el context real, si anés cap amunt (és a
        dir, si el benefici cresqués sense límit en apujar el preu)?
  \item \textbf{(Síntesi del cas)} Redacta en dues o tres frases la recomanació final
        que faries al grup de joves: a quin preu vendre, quin benefici n'esperen i
        entre quins preus, com a mínim, no perdran diners. Justifica-ho amb els valors
        que has calculat.
\end{exlist}
\end{exercicis}

% ---------------------------------------------------------------------
\fullsolucions{Sessió 8}
\begin{solucions}
\begin{sollist}
  \item $B(10)=-500+1500-400=600\eur$ i $B(20)=-2000+3000-400=600\eur$: coincideixen.
        Té sentit perquè $10$ i $20$ són simètrics respecte al vèrtex $x=15$ (a la
        mateixa distància, $5$ a cada costat), i la paràbola és simètrica.
  \item (a) $-2x^2+20x-32=0 \Rightarrow x^2-10x+16=0 \Rightarrow (x-2)(x-8)=0$, és a
        dir $x=2\eur$ i $x=8\eur$. \quad (b) $x_V=\dfrac{-20}{2\cdot(-2)}=5$;
        $B(5)=-50+100-32=18\eur$. Preu òptim $5\eur$, benefici màxim $18\eur$.
  \item $B(8)=-320+1200-400=480\eur$ (a l'esquerra del màxim: convé \textbf{apujar} el
        preu). $B(22)=-2420+3300-400=480\eur$ (a la dreta del màxim: convé
        \textbf{abaixar} el preu). Tots dos donen $480\eur$ per simetria.
  \item $B(25)=-3125+3750-400=225\eur$. És força menys que els $725\eur$ del preu
        òptim ($15\eur$): a partir del vèrtex, apujar més el preu \emph{redueix} el
        benefici, perquè es venen menys samarretes. «Com més car, més es guanya» és
        fals passat el punt òptim.
  \item Va cap avall perquè el terme quadràtic és negatiu ($a=-5<0$): en apujar molt el
        preu, la gent compra menys i el benefici cau. Si anés cap amunt, voldria dir que
        el benefici creixeria indefinidament com més car fos, cosa que no té sentit real:
        ningú no compraria a preus altíssims.
  \item Resposta oberta. Ha de recollir: preu recomanat $15\eur$, benefici esperat
        $\approx 725\eur$, i marge sense pèrdues entre $\approx 3\eur$ i $\approx 27\eur$,
        justificat amb el vèrtex i els punts de tall.
\end{sollist}
\end{solucions}
